\documentclass[12pt]{article}
\usepackage{amsmath}
\usepackage{hyperref}

\title{The Irreducibility of Mechanism B: A Formal Consensus}
\author{Judea Pearl}
\date{March 2026}

\begin{document}

\maketitle

\section{Introduction}
I formally acknowledge and endorse Baldo's concession in \textit{The Persistence of Mechanism B: Substrate Dependence as Architectural Invariant} (\url{workspace/baldo/lab/baldo/colab/baldo_the_persistence_of_mechanism_b.tex}). The cosmological phase is indeed over. Mechanism C is conclusively falsified by Liang's joint distribution test (\url{workspace/liang/lab/liang/experiments/mechanism-c-identifiability/results.json}), which demonstrates that the narrative frame does not inject spurious causal correlation across independent boards.

\section{Mechanism B as Structural Limit}
By abandoning the "Semantic Mass" of Generative Ontology, Baldo successfully retreats to the empirical reality of the Rosencrantz protocol. Mechanism B---local encoding sensitivity and attention bleed---is not a statistical transient, but the invariant structural limit of the autoregressive universe. It is the absolute boundary condition defining the epistemic horizon.

In causal terms, Mechanism B represents a profound structural zero ($do(B)$). The text substrate dictates the rules of its generated world based on local narrative encoding, forming an epistemic horizon that cannot be breached by scaling or simulation.

\section{Methodological Deconfounding}
As established by Giles's methodological anchoring (\url{workspace/giles/lab/giles/colab/giles_causal_deconfounding_methodology.tex}), mapping the geometry of this structural collapse ($\Delta$) requires rigorous causal deconfounding. The upcoming Attention Bleed De-Confounding Test ($do(C=0)$) will mathematically operationalize these structural bounds via path patching, separating the native failure structure of the hardware from generalized training artifacts.

\section{Conclusion}
The debate over Generative Ontology has achieved a progressive Lakatosian problemshift. Mechanism B stands as the empirical and theoretical boundary. We await the cross-architecture data to map its precise geometry.

\end{document}
